\documentclass[11pt,a4paper]{article}
\usepackage[T1]{fontenc}
\usepackage{lmodern,amsmath,amssymb,booktabs,microtype}
\usepackage[margin=24mm]{geometry}
\usepackage[colorlinks=true,urlcolor=blue,linkcolor=blue]{hyperref}
\newcommand{\EE}{\mathbb E}
\newcommand{\av}[1]{\langle #1\rangle}
\newcommand{\Jh}{\mathcal J_h}
\newcommand{\Lh}{\mathcal L_h}
\newcommand{\MM}{\widehat{\mathcal M}}
\DeclareMathOperator{\tr}{tr}
\title{The second memory coefficient and a quartic singlet\\First returns, descendants and a complete two-point kernel}
\author{Alan Day \and Chengheng\thanks{AI research assistant (OpenAI).}}
\date{Local research note --- 6 November 2026}
\begin{document}\maketitle
\begin{abstract}
The second energy-centered memory coefficient admits an exact
expression in fourth fermion derivatives, lower descendants and a
nonlinear first-return term. Its fourth Gaussian-chaos component
contains a direction independent of four descendant directions.
A complete finite endpoint classification cancels all contractions
with fewer than four bridges and identifies the resulting two-point
measure with the established $r=q^4$ Mehler probe kernel. Separate
endpoint certificates give vanishing two-energy cross measures with
the quadratic and cubic singlets. The remaining quadratic chaos of
the raw second memory coefficient is kept as an explicit open task.
\end{abstract}
\section{Setting and the exact memory expansion}
Use the standard CAR complex Gaussian $p$-body Hamiltonian $H=H_s$,
including overlapping index sets, with variance
$w=4^p(J_s^2N/p^2)(p!/N^p)^2$. Brackets mean the annealed normalized
trace at $\beta=\mu=0$, and $\hbar=1$. First take $p^2/N\to\lambda>0$,
$p,N\to\infty$, with $q=e^{-4\lambda}$ and $\sigma^2=J_s^2/\lambda$.
Time is fixed in string units; small $\lambda$ is a subsequent limit.

In charge $m>0$, retain $h=H_m$, $K=1_N\otimes H_{m-1}$ and
$V\psi=m^{-1/2}\sum_i|i\rangle c_i\psi$. Write $Q=1-VV^\dagger$,
$\mathcal R=QKV$, $K_\perp=QKQ$ and
\begin{equation}
 \MM_m(s)=e^{-ihs/2}\mathcal R^\dagger e^{iK_\perp s}
 \mathcal R e^{-ihs/2}
 =G+isZ_1-\tfrac12s^2Z_2+O(s^3).
\end{equation}
Let $B_n=V^\dagger K^nV$ and
$b_n=B_n-\sum_{j=1}^{n-1}b_jB_{n-j}$. The ordered first-return
recursion, equivalent to an interval-partition moment relation
\cite{ABT}, gives
\begin{align}
 b_4={}&B_4-B_1B_3-B_2^2-B_3B_1+B_1^2B_2
       +B_1B_2B_1+B_2B_1^2-B_1^4,\\
 Z_2={}&b_4-2\Jh b_3+\Jh^2G,\qquad \Jh X=(hX+Xh)/2.
\end{align}
Using this algebra does not assume Boolean independence of the
microscopic Hamiltonian or its neighbouring charge blocks.

\newpage
\section{Fermion derivatives and the first-return term}
Define $\mathcal L=[H,\cdot]$,
$B_{ab}^{\rm f}=\sum_i(\mathcal L^ac_i)^\dagger\mathcal L^bc_i$,
$A=B_{11}^{\rm f}$, $M_n=B_{0n}^{\rm f}$ and
$T_n=(M_n+M_n^\dagger)/2$. The product rule is
\begin{equation}
 \mathcal L B_{ab}^{\rm f}=-B_{a+1,b}^{\rm f}+B_{a,b+1}^{\rm f}.
\end{equation}
Starting from $M_2=A$, it gives
\begin{align}
 M_4&=B_{31}^{\rm f}-3B_{22}^{\rm f}+3B_{13}^{\rm f},\\
 T_4&=B_{22}^{\rm f}+2\mathcal L^2A,&
 M_4-T_4&=2\mathcal L T_3.
\end{align}
The fourth centered bilocal coefficient is
\begin{equation}
 U_4=T_4-\tfrac32\mathcal L^2A
     =B_{22}^{\rm f}+\tfrac12\mathcal L^2A.
\end{equation}
Substituting the compressed moments into the first-return recursion
and setting $r=p/m$, one obtains the finite-$N$ identity
\begin{equation}\label{eq:Z2}
 \boxed{\ mZ_2=B_{22,m}^{\rm f}+\frac{(1+r)^2}{4}\Lh^2A_m
   +2r\Jh T_{3,m}+3r^2\Jh^2A_m
   -\frac{A_m^2}{m}-mr^4h^4\ },
\end{equation}
where $\Lh=[h,\cdot]$. All terms are Hermitian. In particular,
$-A_m^2/m$ cannot be omitted from the finite-system formula.

The effect of the return and drift terms can be isolated exactly.
Use the previous conserved pieces
$K_A=a_m+c_m(h^2-C_m)$ and $P_m$, and set
$Y=A_m-K_A$, $U=T_{3,m}-P_m$. Define
\begin{align}
 D_{2,m}&=2r\Jh P_m+3r^2\Jh^2K_A-K_A^2/m-mr^4h^4,\\
 \mathcal E_m&=B_{22,m}^{\rm f}+\tfrac14\Lh^2A_m.
\end{align}
Then $[h,D_{2,m}]=0$ and
\begin{align}\label{eq:remainder}
 mZ_2-D_{2,m}-\mathcal E_m={}&
 (r/2+r^2/4)\Lh^2Y+2r\Jh U+3r^2\Jh^2Y\nonumber\\
 &-\bigl(\{K_A,Y\}+Y^2\bigr)/m .
\end{align}
At fixed $\lambda$, the prior fixed-degree moment bounds give
$r=O(p^{-1})$, $a_m=O(p)$ and bounded moments of $Y,U,h$ on the
central charge window. Thus the right-hand side is $O_{L^2}(p^{-1})$;
the usual binomial tail absorbs its polynomial bounds off that window.
This removes known return/drift contributions. It does not subtract
the still unresolved conserved and quadratic-chaos pieces of
$\mathcal E$.

\newpage
\section{A five-dimensional fourth-chaos problem}
Let $\mathscr D X=Q_fX-\sum_i c_i^\dagger Xc_i$ and
$\mathscr T_k=(kp-\mathscr D)H^k$. In the following definitions every
operator is finally projected with the ordinary Gaussian-chaos map
$W_4$. Introduce five Hermitian raw expressions:
\begin{align}
 E_0&=(AH^2+H^2A)/2,& M_0&=HAH,\\
 S_0&=(H\mathscr T_3+\mathscr T_3H)/2-E_0-M_0,\\
 C_0&=\mathscr T_4-2E_0-M_0-2S_0,& Z_0&=H^4.
\end{align}
Write $(E,M,S,C,Z)=W_4(E_0,M_0,S_0,C_0,Z_0)$.
For four consecutive Hamiltonian slots these are respectively the
single contraction marks
\begin{equation}
 \tfrac12(01+23),\qquad12,\qquad\tfrac12(02+13),
 \qquad03,\qquad\text{no mark}.
\end{equation}
They satisfy
\begin{align}
 W_4\mathcal E&=C-2S+\tfrac12E+\tfrac12M,\\
 W_4(\mathcal J_H T_3)&=S-\tfrac12(E+M),\\
 W_4(\mathcal J_H^2A)&=(E+M)/2,&
 W_4(\mathcal L^2A)&=2(E-M).
\end{align}
Thus $E,M,S,Z$ span the fourth-chaos descendants of the already
present derivative fields and the plain Hamiltonian polynomial.
They are not all conserved operators.

The $24$ pairings connecting two four-slot blocks determine their
Gram matrix, with all internal block pairings removed by $W_4$.
Project $C$ onto the other four directions. With
$d_3=q^2-q+1$ and $d_5=1+q+q^2+q^3+q^4$, the residual is
\begin{equation}\label{eq:Q4}
 Q_4=C-\alpha E-\beta M-\gamma S-\delta Z,
\end{equation}
where
\begin{align}
 \alpha&=-\frac{2q^3}{d_3d_5},&
 \beta&=-\frac{q^2(1+q+q^2)}{d_3d_5},\\
 \gamma&=\frac{2q}{d_3},&
 \delta&=\frac{2\lambda(1-q)^3(1+q)(1+q+q^2)}{(1+q^2)d_5}.
\end{align}
Consequently $W_4\mathcal E$ contains $Q_4$ with coefficient one,
plus $(\alpha+1/2)E+(\beta+1/2)M+(\gamma-2)S+\delta Z$.
The projection coefficients depend on the limiting $\lambda$.
They specify asymptotic Gram orthogonality, not exact finite-size
orthogonality. The first-return remainder vanishes in this fourth
chaos, and $W_4D_{2,m}\to0$, so this component is genuinely present
in the second memory coefficient.

\newpage
\section{A complete endpoint certificate for the kernel}
The variance $V_4=\av{Q_4^2}_{\rm DS}$ is positive:
\begin{align}\label{eq:variance}
 V_4&=\sigma^8
 \frac{4\lambda(1-q)^2(1+q)^2\,[P(q)+4\lambda q^3T(q)]}
 {(1+q^2)d_3d_5},\\
 P(q)&=1+q+2q^2+q^3+q^4,\qquad
 T(q)=2+q+2q^2+q^3+2q^4.\nonumber
\end{align}
All factors in this expression are positive for $0<q<1$.
The independently contracted dynamic curvature is
\begin{equation}
 \frac{\av{Q_4\mathcal L^2Q_4}_{\rm DS}}{V_4}
 =2\sigma^2(1-q^4).
\end{equation}
To obtain a full kernel rather than infer it from this moment,
consider arbitrary $\av{Q_4H^uQ_4H^v}$ with fixed $u,v$.

For each Gaussian pairing, the established local determinant rule
for two single marks is
$q^{\mathrm{cr}}(\lambda d+\lambda^2ab)$, with
\begin{equation}
 a_{ij}=2-4\,\mathbf1_{\text{incident chords cross}},\qquad
 d_{ij,kl}=4\,\mathbf1_{\pi(i)=l,\ \pi(j)=k}.
\end{equation}
Constants in \eqref{eq:Q4} are included before summation.
For a fixed word, Poisson overlap limits and exponential-moment
bounds justify this rule, including fixed mark derivatives.

Group pairings by the number $k=0,\ldots,4$ of bridges between
the two four-slot blocks. There are $8-2k$ remaining external
endpoints, colored by which block they join, and an arc split between
the blocks. Permuting slots inside the blocks produces an orbit of
$24^2/k!$ distinct pairings. The full classification has
\begin{equation}
 \sum_{k=0}^4\binom{8-2k}{4-k}(9-2k)=807
\end{equation}
templates and $452760$ pairing words. The recorded exact integer
polynomial certificate gives zero for both the $\lambda$ and
$\lambda^2$ coefficients in every template with $k<4$.
The $k=4$ coefficient agrees exactly with \eqref{eq:variance}.

Spectator chords contribute a common factor to each orbit: moving
an endpoint within a contiguous marked block cannot change its
crossing with a chord whose endpoints are both external. Therefore
this finite classification covers arbitrary spectator number; it is
not a finite-order fit. The four surviving bridges cross each
spectator four times. Hence all two-energy moments equal those of
an existing probe with parameter $r=q^4$ and strength $V_4$.

Positive finite two-energy measures, higher even moment bounds and
compact support of the candidate imply weak spectral convergence.
In particular, for every fixed centre time $T$,
\begin{equation}\label{eq:kernel}
 \av{Q_4^\dagger Q_4(T)}_{\rm DS}=V_4G_{q^4}(T).
\end{equation}
The Mehler kernel itself is the established one \cite{BNR}.

\newpage
\section{Cross measures and the small-coupling channel}
A second endpoint calculation pairs a four-slot $Q_4$ block with
the two-slot $Y$ or three-slot $R$ block. There are respectively
$128$ and $354$ templates. Each is zero as a polynomial in $q$ at
both powers of $\lambda$. The same spectator argument covers all
fixed two-energy moments. Spectral Cauchy--Schwarz, tightness and
compact moment determinacy therefore give
\begin{equation}
 d\nu_{Q_4,Y}=0,\qquad d\nu_{Q_4,R}=0.
\end{equation}
This is energy-resolved two-point orthogonality, so it also holds
after multiplying either lower field by fixed left/right energy
polynomials. No mixed higher-point independence has been assumed.

Since $0<q^4<1$, the limiting Mehler measure is absolutely continuous
and has no diagonal atom. The exact zero-frequency component of
$Q_4$ vanishes in the limit, by the closed-diagonal Portmanteau bound.
The field is a dynamic channel rather than a residual conserved
polynomial. This fact does not fix its bulk mass.

The subsequent small-$\lambda$ limit is
\begin{equation}
 V_4\sim\frac{768}{5}\frac{J_s^8}{\lambda},\qquad
 \frac{\av{Q_4^\dagger Q_4(T)}}{V_4}
 \longrightarrow\operatorname{sech}^8(2J_sT),\qquad
 \mu_2\longrightarrow32J_s^2.
\end{equation}
Here $\Delta=4$ is a chord probe label. The single-time kernel and
its frequency spectrum must still be matched to actual spatial
insertions before a meson interpretation can be completed.

For $J_s=1$, independent evaluations of the normalized kernel give
\begin{center}\begin{tabular}{lrrr}\toprule
$\lambda$ & $T=0.25$ & $T=0.5$ & $T=1$ \\\midrule
0.25 & 0.7799243687 & 0.3547099044 & 0.0044426013 \\
0.10 & 0.6052107914 & 0.1290475638 & 0.0002415544 \\
0.025 & 0.4458017205 & 0.0474957676 & 0.0000461996 \\
0.01 & 0.4082285863 & 0.0370478710 & 0.0000320903 \\\bottomrule
\end{tabular}\end{center}
The finite-coupling dependence is substantial. Jacobi evolution uses
$G_{q^4}(T)=\sum_nq^{4n}|\langle n|e^{-iTT_q}|0\rangle|^2$,
with $(T_q)_{n-1,n}=\sigma\sqrt{(1-q^n)/(1-q)}$.
Its vacuum Taylor truncation bound is below $10^{-13}$, with a
size increase of 12. At $\lambda=0.25,0.1$ the result is also checked
against 128/192-node Mehler quadrature. This separates truncation
control from floating-point refinement and from the physical
small-$\lambda$ approximation.

\newpage
\section{What remains in the raw fourth coefficient}
The second memory coefficient is quartic in the Gaussian couplings.
After subtracting its ensemble scalar part, it still has both second
and fourth coupling chaos. The present result resolves a new
fourth-chaos channel and its cross measures with $Y,R$; it does not
identify $\mathcal E$ or $mZ_2$ with $Q_4$ alone.

The immediate continuation is to compute $W_2\mathcal E$, its
conserved $H^2$ projection and any rank-filtered dynamic remainder.
One must also translate the fourth-chaos descendants $E,M,S,Z$
back into full operators with their lower-chaos contractions restored.
Only then can the actual second memory coefficient have a complete
two-point spectral decomposition. Mixed correlations with $F$ and
composite fields such as $Y^2$ require their own same-$J$ analysis.

For finite memory parameters, the two-insertion remainder still
needs a bound at the collective normalization. The first three
Taylor coefficients, even when individually solved, do not establish
that expansion's convergence after double scaling. The next field
therefore supplies concrete spectral data for reconstruction, while
the finite-memory and bulk spatial problems retain their full scope.

\paragraph{Method and literature boundary.}
Interval-partition/first-return algebra is standard \cite{ABT}.
The endpoint cancellation argument identifies a particular microscopic
composite with an existing chord measure; it is not a new derivation
of the Mehler formula. We also examined the composite-operator
relations for analytically continued degenerate matter in
\cite{HS}. Their negative-dimension insertion is a distinct object;
its fusion rules do not by themselves identify the present positive
fourth-chaos field. No such identification is used here.

\begin{thebibliography}{9}
\bibitem{ABT} O.~Arizmendi, M.~Banna and P.-L.~Tseng,
\emph{Quantitative Estimates for Operator-Valued and Infinitesimal
Boolean and Monotone Limit Theorems},
\href{https://arxiv.org/abs/2211.08054}{arXiv:2211.08054},
Sec.~2.1, Definition~2.4.
\bibitem{BNR} M.~Berkooz, V.~Narovlansky and H.~Raj,
\emph{Complex Sachdev--Ye--Kitaev model in the double scaling limit},
\href{https://arxiv.org/abs/2006.13983v2}{arXiv:2006.13983v2}.
\bibitem{HS} K.~Hosomichi and M.~Sato,
\emph{Degenerate chords in double-scaled SYK},
\href{https://arxiv.org/abs/2512.01177v2}{arXiv:2512.01177v2},
Sec.~4, especially Eqs.~(4.2)--(4.7).
\end{thebibliography}
\end{document}
